% New separate result after S12--S15; 2026-11-22 UTC.
% GPT-6 Astra Pro, /root/fixed_surface; effort not exposed.

\section[HD-F1-S16: the ordinary simple norm sheaf]{HD-F1-S16: the actual norm object is an ordinary simple sheaf}
\label{sec:f1-norm-fourfold-derived}

Restrict the actual cover parameter space of HD-F1-S13 to its nonempty
open intersection with the open in Schwarz's Theorem 1.1 excluding
$g^2_8$ on the covering curve. For each of the three components let
$a:T\to Y\subset P$ be the actual Abel morphism, and let
$\nu:Y^\nu\to Y$ be the finite normalization. Then
\[
 R^qa_*\mathcal O_T=0\quad(q>0),\qquad
 Ra_*\mathcal O_T=\nu_*\mathcal O_{Y^\nu},\qquad
 \operatorname{Hom}_P(Ra_*\mathcal O_T,Ra_*\mathcal O_T)=\mathbb C,
 \tag{F1S16.1}
\]
where the closed inclusion of $Y$ in $P$ is understood in the
pushforward. In particular the specified actual candidate is an
ordinary simple coherent sheaf and
\[
 \operatorname{Ext}^k_P(Ra_*\mathcal O_T,Ra_*\mathcal O_T)=0
                           \quad(k<0).
 \tag{F1S16.2}
\]
Its full character is exactly HD-F1-S14; no replacement of the object
or change of its middle determinant component occurs. The normalization
$Y^\nu$ has rational singularities. The original image $Y$ remains the
actual singular image of S15; no normality assertion for $Y$ is needed.

\paragraph{1. The scheme fibres and their actual normal maps.}
The scheme fibre of the Abel map $\operatorname{Sym}^8Z\to
\operatorname{Pic}^8Z$ at $L$ is the complete projective space
$\mathbb P H^0(Z,L)$. This is the fine divisor description: effective
Cartier divisors with associated line bundle $L$ are nonzero sections
of $L$ modulo scalars, and this description holds for families over
arbitrary test schemes. Since $\operatorname{Nm}L=K_C$, all such
sections belong to the norm fibre. The connected projective fibre lies
in a single $T_j$. Thus the scheme fibre of the restricted $a$ is the
same complete projective space, not merely its set of closed points.
By Schwarz's theorem every nonempty such fibre is a reduced point or
a reduced complete pencil $E\simeq\mathbb P^1$.

At every divisor $D\in E$, the divisor exact sequence identifies the
kernel of the full Abel differential as
$H^0(Z,L)/H^0(Z,\mathcal O_Z)$, where the inclusion of constants is
multiplication by the section defining $D$. It is one dimensional
and is exactly $T_DE$. Restriction to $T$ therefore gives the
fiberwise injective map
\[
          N_{E/T}\longrightarrow T_{P,a(E)}\otimes\mathcal O_E
                                      \simeq\mathcal O_E^{\oplus8}.
 \tag{F1S16.3}
\]
Its rank is three at every point. A map of vector bundles with constant
rank and injective fibres is a subbundle, so the dual map is an actual
surjection $\mathcal O_E^{\oplus8}\twoheadrightarrow N^*_{E/T}$.
Consequently
\[
 N^*_{E/T}\simeq\mathcal O_{\mathbb P^1}(b_1)\oplus
 \mathcal O_{\mathbb P^1}(b_2)\oplus\mathcal O_{\mathbb P^1}(b_3),
                  \qquad b_i\geq0.
 \tag{F1S16.4}
\]
Here splitting is the usual vector-bundle splitting on the projective
line; the nonnegative degrees follow directly from global generation.
In fact $b_1+b_2+b_3=3$: by adjunction,
$\deg N_{E/T}=\deg K_E-\deg K_T|_E=-2-1=-3$, because
$\vartheta|_E=0$ and $\eta|_E=\mathcal O_E(1)$.
We do not assume a specific splitting type. All that is needed below is
\[
       H^q(E,\operatorname{Sym}^m N^*_{E/T})=0
                    \quad(m\geq0,\ q>0),
 \tag{F1S16.5}
\]
which follows by expanding the symmetric power of (F1S16.4) into
nonnegative-degree line bundles and using the cohomology of
$\mathcal O_{\mathbb P^1}(b)$ for $b\geq0$.

\paragraph{2. Formal functions gives every higher direct-image vanishing.}
Fix $y\in Y$ with fibre $E=\mathbb P^1$. Since the scheme fibre is this
reduced pencil, its ideal in $T$ is exactly
$\mathcal I=\mathfrak m_y\mathcal O_T$. Both $T$ and $E$ are smooth,
so this is a regular embedding and
\[
 \mathcal I^m/\mathcal I^{m+1}
                    \simeq\operatorname{Sym}^m N^*_{E/T}.
 \tag{F1S16.6}
\]
Moreover $(\mathfrak m_y\mathcal O_T)^m=
\mathfrak m_y^m\mathcal O_T$. Hence the infinitesimal fibres used by
the theorem on formal functions are precisely the thickenings
$E_m$ defined by $\mathcal I^{m+1}$; no comparison between different
adic topologies is required. Starting with $H^{>0}(E,\mathcal O_E)=0$,
use the actual short exact sequences
\[
 0\longrightarrow\operatorname{Sym}^mN^*_{E/T}
 \longrightarrow\mathcal O_{E_m}
 \longrightarrow\mathcal O_{E_{m-1}}\longrightarrow0
\]
and (F1S16.5) to prove $H^q(E_m,\mathcal O_{E_m})=0$ for every $q>0$
and every $m$. At a reduced-point fibre, every infinitesimal fibre has
zero-dimensional support, so the same higher cohomology vanishing
holds for every thickening. At a point outside $Y$ the fibre is empty.
The theorem on formal functions (Stacks Project, Theorem 30.20.5,
Tag 02OC, and its stalk formulation Lemma 30.20.7) now gives the zero completion of the stalk of
$R^qa_*\mathcal O_T$ at every closed point for every $q>0$.
These sheaves are coherent by properness; completion of a finitely
generated module over a noetherian local ring is faithful. Thus all
these higher direct images vanish. This proves the first equality of
(F1S16.1), in every positive degree.

\paragraph{3. The normalization, rational singularities and simplicity.}
The smooth $T$ is normal and $a$ is birational, so its dominant map to
$Y$ factors through the normalization:
\[
          T\xrightarrow{b}Y^\nu\xrightarrow{\nu}Y.
 \tag{F1S16.7}
\]
The morphism $b$ is proper birational with normal target. The usual
normal-target form of Stein factorization (equivalently the
regular-function form of Zariski's Main Theorem) gives
$b_*\mathcal O_T=\mathcal O_{Y^\nu}$. One can check it in the common
function field: a regular function upstairs has no pole over any
codimension-one point of the normal target, so normality makes it a
regular function on that target. Since $\nu$ is finite, its pushforward
is exact and faithful. Thus the first part of the proof also gives
$R^qb_*\mathcal O_T=0$ for $q>0$. This is precisely the rational
singularity property of $Y^\nu$ for the displayed resolution $b$.
It also proves the second equality in (F1S16.1).

Finally let $M=\nu_*\mathcal O_{Y^\nu}$, considered as a sheaf on $Y$.
It is a torsion-free rank-one $\mathcal O_Y$-module inside the common
function field. An $\mathcal O_Y$-linear endomorphism is multiplication
by an element of that field at the generic point, and this determines
it everywhere by torsion-freeness. Preserving $M$ forces this element
to belong to $M$, by evaluating the endomorphism at $1$; conversely
every element of $M$ gives multiplication preserving $M$. Hence
\[
       \mathcal H om_Y(M,M)=M,\qquad
       \operatorname{Hom}_P(i_*M,i_*M)
          =H^0(Y^\nu,\mathcal O_{Y^\nu})=\mathbb C.
\]
The last equality uses that $Y^\nu$ is integral and projective.
The degree-negative Ext vanishing (F1S16.2) is now the actual
orthogonality of the standard derived-category $t$-structure for this
ordinary sheaf. No truncation or direct-summand substitution is used.

\paragraph{The precise remaining trace gate.}
The same perfect object now has its full prescribed character,
simplicity, and negative self-Ext vanishing established. The next map
is its actual trace
\[
 \sigma_{\mathcal E}:\operatorname{Ext}^2_P(\mathcal E,\mathcal E)
                                      \longrightarrow HH_{-2}(P),
 \qquad\mathcal E=i_*\nu_*\mathcal O_{Y^\nu}.
 \tag{F1S16.8}
\]
The formal-fibre argument proves no rank assertion for this map.
Its kernel, and any specified finite-group invariant part of that
kernel, remain to be computed for the object just identified.
